\documentclass[11pt,a4paper]{article}
\usepackage{amsmath,amssymb,amsfonts}
\usepackage{geometry}
\usepackage{hyperref}
\geometry{margin=1in}

\title{Logarithmic quantum corrections to the Schwarzschild-AdS metric and modified holographic entanglement entropy}
\author{Discovery AI Autonomous Engine}
\date{October 16, 2026}

\begin{document}
\maketitle

\begin{abstract}
The holographic principle and AdS/CFT correspondence provide a powerful framework for understanding quantum gravity and the emergence of spacetime from a boundary conformal field theory (CFT). In this manuscript, we explore the implications of logarithmic quantum corrections to the Schwarzschild-AdS metric, which arise from the renormalization group flow of the quantum stress-energy tensor in the boundary CFT. These corrections lead to a running of the AdS/CFT coupling constant, governed by the beta function determined by the renormalization group equations for the boundary CFT. We derive the equations for the running of the AdS/CFT coupling constant and discuss the asymptotic behavior of the logarithmic quantum corrections. We also propose testable predictions for the scale-dependent modification of holographic entanglement entropy. Our results suggest that logarithmic quantum corrections play a crucial role in understanding the quantum properties of black holes in anti-de Sitter spacetimes and may have important implications for the information paradox and the emergence of spacetime.
\end{abstract}

\section{Introduction and Theoretical Motivation}
The holographic principle and AdS/CFT correspondence have revolutionized our understanding of quantum gravity and the emergence of spacetime from a boundary conformal field theory (CFT). The holographic principle states that the description of a volume of space can be thought of as the edge of that space, with the theory residing on the boundary. The AdS/CFT correspondence is a specific realization of the holographic principle, where a quantum gravity theory in a certain anti-de Sitter (AdS) spacetime is equivalent to a CFT on the boundary of that spacetime.

In this manuscript, we focus on the implications